\chapter{Multilinear Kakeya and the hairbrush bound}
\label{chap:multilinear}

Two older estimates enter the structural argument. Wolff's hairbrush
bound prevents the critical exponent from reaching $1$, while
multilinear Kakeya forces a nearly extremal family to be almost planar
at a typical shaded point. We give the statements needed for these
uses and the numbered ingredients of the free manuscript of Carbery
and Valdimarsson~\cite{CV}. Their argument works in arbitrary dimension;
the completed Lean interfaces used here specialize to dimension three.

\section{The endpoint multilinear estimate}

\begin{theorem}[Multilinear Kakeya; CV, Theorem 1]
\label{thm:cv-general}
Let $2\leq d\leq n$, and let $\T_1,\ldots,\T_d$ be finite indexed
families of unit neighbourhoods of complete lines in $\R^n$. Write
$e(T)$ for a unit direction and
$v_1\wedge\cdots\wedge v_d$ for the unsigned volume of the
parallelepiped spanned by these vectors. For nonnegative weights $a_T$,
\[
 \int_{\R^n}
 \left(\sum_{T_j\in\T_j}
     \prod_{j=1}^d a_{T_j}\one_{T_j}(x)
     \;e(T_1)\wedge\cdots\wedge e(T_d)\right)^{1/(d-1)}dx
 \leq C_{n,d}\prod_{j=1}^d
                  \left(\sum_{T\in\T_j}a_T\right)^{1/(d-1)}.
\]
\end{theorem}
\begin{proof}
\uses{prop:cv-reduction,thm:cv-mollified-visibility,lem:cv-mollified}
Apply polynomial visibility to a normalized weight on the unit cubes.
Directional surface areas give a factorization of the wedge weight;
the cylinder estimate bounds each column sum. The finite reduction
then gives the integral inequality. Integer multiplicities, homogeneity,
and approximation recover general nonnegative coefficients.
\end{proof}

\begin{theorem}[The shaded three-dimensional estimate; WZ1, Theorem 4.2]
\label{thm:multilinear-kakeya}
\lean{Kakeya.CV.MultilinearKakeyaThreeStatement,Kakeya.CV.multilinear_kakeya_three}
\lean{Kakeya.Assouad.paper_cv_broad_set_estimate}
\leanok
\uses{def:tube,def:shading}
For $0<\delta\leq1$, any finite indexed family $\T$ of
$\delta$-tubes and any measurable shading satisfy
\[
 \int_{\R^3}
  \left(\sum_{T_1,T_2,T_3\in\T}
    \prod_{j=1}^3\one_{Y(T_j)}(x)
    |\det(e(T_1),e(T_2),e(T_3))|\right)^{1/2}dx
       \lesssim (\delta^2\#\T)^{3/2}.
\]
\end{theorem}
\begin{proof}
\uses{thm:cv-general}
\leanok
Take three copies of the same family, enlarge each finite tube to
the neighbourhood of its complete axis, and rescale by $\delta$.
Restriction to the shading decreases the integrand. The Jacobian
contributes $\delta^3$, which is exactly the power on the right.
The sticky proof reaches the unit-scale theorem through
\texttt{paper\_cv\_broad\_set\_estimate}, using the paper's cropped
full-line shadings and a rescaling by $6\delta$. The linked
\texttt{multilinear\_kakeya\_three} theorem is a separately proved
finite-segment formulation; its proof term is not used by the sticky endpoint.
\end{proof}

\begin{proposition}[Finite factorization; CV, Proposition 2]
\label{prop:cv-reduction}
\lean{Kakeya.CV.multilinear_kakeya_preliminary_reduction}
It suffices for Theorem~\ref{thm:cv-general} to construct, for every
finitely supported $M\geq0$ with $\sum_Q M(Q)^n=1$, nonnegative
numbers $S_j(Q,T)$ such that
\[
 e(T_1)\wedge\cdots\wedge e(T_d)\,M(Q)^n
       \lesssim\prod_{j=1}^d S_j(Q,T_j),
 \qquad
 \sum_{Q:Q\cap T\ne\varnothing}S_j(Q,T)\lesssim1,
\]
where the first inequality is required when all the tubes meet $Q$.
The linked finite inequality is its three-dimensional specialization.
\end{proposition}
\begin{proof}
Let $F(Q)$ be the sum of the wedge weights over incident tuples.
Choose $M(Q)^n$ proportional to $F(Q)^{1/(d-1)}$, first on a finite
set of cubes. Sum the factorization and apply H\"older's inequality
and the column bounds. Rearrangement gives the desired bound for
$\sum_Q F(Q)^{1/(d-1)}$; exhaust the lattice by finite subsets.
\end{proof}

\section{Visibility and polynomial surfaces}

\begin{definition}[Directional area and visibility; CV, Section 3]
\label{def:visibility}
For a hypersurface $Z$ with unit normal $n_Z$ almost everywhere, put
\begin{align*}
 \operatorname{surf}_e(Z)&=\int_Z|e\cdot n_Z|\,d\cH^{n-1},\\
 K(Z)&=\left\{u\in B(0,1):
                    \int_Z|u\cdot n_Z|\,d\cH^{n-1}\leq1\right\},\\
 \operatorname{vis}(Z)&=|K(Z)|^{-1/n}.
\end{align*}
The body $K(Z)$ is convex and centrally symmetric. For a polynomial
$p$, write $Z_p=\{p=0\}$. A mollified directional area averages the
area functional over nearby polynomial parameters; its associated
body and visibility are denoted by $K_\tau$ and
$\operatorname{vis}_\tau$.
\end{definition}

\begin{lemma}[Cylinder estimate; CV, Lemma 1]
\label{lem:cv-cylinder}
\lean{Kakeya.CV.polynomial_cylinder_estimate}
\uses{def:visibility}
If $T$ is a unit neighbourhood of a line and $p\ne0$ has degree at
most $k$, then
\[
                \operatorname{surf}_{e(T)}(Z_p\cap T)\lesssim_n k.
\]
\end{lemma}
\begin{proof}
Project onto the hyperplane perpendicular to $e(T)$. Almost every
fibre meets the zero set at most $k$ times. Integrating this count
over a bounded cross-section proves the estimate.
\end{proof}

\begin{lemma}[Visibility from directional areas; CV, Lemma 2]
\label{lem:cv-directions}
\lean{Kakeya.CV.visibility_directional_surface_bound}
\uses{def:visibility}
Suppose $1\lesssim\operatorname{surf}_e(Z)\lesssim D$ for every
unit $e$. For unit vectors $v_1,\ldots,v_d$, $1\leq d\leq n$,
\[
 (v_1\wedge\cdots\wedge v_d)^{1/n}\operatorname{vis}(Z)
   \lesssim D^{(n-d)/n}
              \prod_{j=1}^d\operatorname{surf}_{v_j}(Z)^{1/n}.
\]
\end{lemma}
\begin{proof}
Complete the vectors to a basis using an orthonormal basis of their
orthogonal complement. The convex hull of the points
$\pm c v_j/\operatorname{surf}_{v_j}(Z)$ lies in $K(Z)$.
Its volume gives the inequality. The Lean link proves this convex
body argument for three vectors in $\R^3$.
\end{proof}

\begin{lemma}[A low-area direction bounds visibility; CV, Lemma 3]
\label{lem:cv-degree}
\lean{Kakeya.CV.visibility_degree_bound}
\uses{def:visibility}
If $Q$ is a unit cube, $p\ne0$, and
$\operatorname{surf}_e(Z_p\cap Q)\lesssim1$ for some unit $e$,
then
\[
             \operatorname{vis}(Z_p\cap Q)^{n/(n-1)}
                          \lesssim 1+\deg p.
\]
\end{lemma}
\begin{proof}
\uses{lem:cv-cylinder}
The John ellipsoid of $K(Z_p\cap Q)$ has one axis bounded below.
Its volume then forces another axis to be short when visibility is
large. The corresponding directional area is large, and the cylinder
estimate bounds it by the degree.
\end{proof}

\begin{remark}[Zero degree and singular polynomial presentations]
\label{rem:cv-polynomial}
CV writes $\deg p$ in Lemma 3. For a nonzero constant polynomial the
zero set is empty, but $K(\varnothing)=B(0,1)$ and visibility is
positive. One must require a positive degree bound $k\geq1$, as in
the Lean interface, or write $1+\deg p$ as above. The later application
uses a positive degree bound. The formal directional area is defined
using a polynomial normal and therefore keeps an explicit negligible
singular-set hypothesis. Squarefree reduction and genericity supply
this hypothesis before the estimate is applied; an arbitrary repeated
factor presentation must not be substituted without that step.
\end{remark}

\begin{lemma}[Mollified estimates; CV, Lemmas 4--6]
\label{lem:cv-mollified}
\uses{lem:cv-cylinder,lem:cv-directions,lem:cv-degree}
For polynomials in a fixed normalized parameter space of degree at
most $k\geq1$, the cylinder, directional visibility, and low-area
visibility bounds remain valid after mollification. Their constants
are independent of the sufficiently small mollification radius.
\end{lemma}
\begin{proof}
Average the cylinder bound. The averaged area remains an even
sublinear function of the direction, so its unit body is convex.
The convex-geometric proofs of the other two estimates apply to this
body without change.
\end{proof}

\begin{theorem}[Polynomial visibility; CV, Theorems 2--3]
\label{thm:cv-mollified-visibility}
\lean{Kakeya.CV.PolynomialVisibilityStatement,Kakeya.CV.polynomial_visibility_complete}
\leanok
\uses{def:visibility}
For finitely many unit cubes in $\R^3$ and nonnegative weights $M(Q)$,
there are a degree bound $k$, a nonzero normalized polynomial of
degree at most $k$, and a mollification radius $\tau>0$ such that
\[
       k\lesssim\left(\sum_Q M(Q)^3\right)^{1/3},\qquad
       \operatorname{vis}_\tau(Z_p\cap Q)\gtrsim M(Q)
       \quad\hbox{for every }Q.
\]
CV Theorem 2 introduces the unmollified assertion in dimension $n$;
Theorem 3 is the mollified version used in the proof. The displayed
statement is the completed finite, three-dimensional version.
\end{theorem}
\begin{proof}
\uses{lem:cv-antipodal,lem:cv-colouring,lem:cv-bisections,
  lem:cv-stability}
\leanok
For each cube, decompose the polynomials with insufficient visibility
by the size and shape of their visibility body. A finite colouring
chooses ellipsoids with a stability margin. Pack each cube with small
translates of the chosen ellipsoid. The bisection estimate leaves
one translate on which a bad polynomial has a strict sign imbalance.
These signs separate each bad set from the closure of its antipode.
The total number of such pairs is controlled by $\sum_Q M(Q)^3$.
Borsuk--Ulam leaves a polynomial outside all the bad sets once the
parameter-space dimension is chosen large enough.
\end{proof}

\begin{proposition}[Prescribing surface area; CV, Proposition 1]
\label{prop:cv-surface-area}
\uses{def:visibility}
If $M$ is finitely supported on the unit cubes of $\R^n$ and each
nonzero $M(Q)$ is at least $1$, there is a nonzero polynomial with
\[
 \deg p\lesssim\left(\sum_Q M(Q)^n\right)^{1/n},\qquad
 \cH^{n-1}(Z_p\cap Q)\gtrsim M(Q)
 \quad(M(Q)>0).
\]
\end{proposition}
\begin{proof}
\uses{thm:borsuk-ulam,lem:cv-isoperimetric}
Subdivide each relevant cube into approximately $M(Q)^n$ smaller
cubes. Borsuk--Ulam gives a polynomial bisecting all of them.
The relative isoperimetric lower bound in each small cube, followed
by summation, gives the asserted surface area.
\end{proof}

\section{The topological construction}

\begin{theorem}[Borsuk--Ulam; CV, Theorem 4]
\label{thm:borsuk-ulam}
If $J\leq N$, every continuous odd map $F:S^N\to\R^J$ has a zero.
\end{theorem}

\begin{lemma}[Antipodal covering; CV, Lemma 7]
\label{lem:cv-antipodal}
\lean{Kakeya.CV.polynomial_bad_set_cover}
\leanok
Let $A_1,\ldots,A_J\subset S^N$, with $J\leq N$. If
$A_i\cap\overline{-A_i}=\varnothing$ for each $i$, the sets
$A_i,-A_i$ do not cover the sphere.
\end{lemma}
\begin{proof}
\uses{thm:borsuk-ulam}
\leanok
Apply Borsuk--Ulam to the odd differences of distance to the antipodal
closures. A common zero cannot belong to any $A_i$ or $-A_i$.
Empty sets can first be omitted.
\end{proof}

\begin{lemma}[A coloured ellipsoid net; CV, Lemma 8]
\label{lem:cv-colouring}
\lean{Kakeya.CV.stable_finite_colour_ellipsoid_net}
There is a dimensional $\alpha>1$ and a family of centred ellipsoids
with finitely many colours such that every centrally symmetric
convex body $K$ is approximated by one of them:
\[
                    \alpha^{-1}K\subset E\subset\alpha K.
\]
The colours can be chosen so that same-colour ellipsoids which are
$\alpha^3$-close must agree. In particular there is at most one
approximating ellipsoid of a given colour.
\end{lemma}
\begin{proof}
Take a separated net in the space of positive definite quadratic
forms, with distance measured by logarithmic homothetic distortion.
Colour its bounded-degree proximity graph. John's theorem then
extends the covering property to centrally symmetric convex bodies.
\end{proof}

\begin{remark}[A margin in the colouring]
\label{rem:cv-colouring}
The stability argument needs separation stronger than mere uniqueness
of the chosen approximation. The formal colouring explicitly retains
the $\alpha^3$ margin. This accounts for the two approximation factors
and the small change of the visibility body when the polynomial varies.
\end{remark}

\begin{lemma}[How many ellipsoids can be bisected; CV, Lemma 9]
\label{lem:cv-bisections}
\uses{def:visibility}
Suppose $p$ lies in a bad visibility band at a cube $Q$, and its
selected ellipsoid $E(p)$ has volume comparable to
$2^{rn}M(Q)^{-n}$. For $0<\eta<1$, the zero set of $p$ can bisect
at most
\[
                     C_n\eta^{-(n-1)}2^{-rn}M(Q)^n
\]
disjoint translates of $\eta E(p)$ inside $Q$. The assertion is used
with a common sufficiently small mollification radius and a fixed
finite collection of possible translates.
\end{lemma}
\begin{proof}
\uses{lem:cv-isoperimetric,lem:cv-mollified}
Each bisection contributes a fixed amount to the sum of the
directional areas weighted by the principal semiaxes. Summing over
disjoint translates and using the directional budgets bounds their
number. Continuity of sign volumes replaces exact bisection by a
fixed approximate bisection for nearby polynomials.
\end{proof}

\begin{lemma}[Stability of the selected ellipsoid; CV, Lemma 10]
\label{lem:cv-stability}
\lean{Kakeya.CV.ellipsoid_selection_stability}
\uses{lem:cv-colouring,def:visibility}
Fix a cube, a visibility band, a colour, and a mollification radius.
If $p_m\to p$ within this class, then $E(p_m)=E(p)$ for all
sufficiently large $m$. Every translate with a fixed label in the
associated packing is eventually unchanged as well.
\end{lemma}
\begin{proof}
Mollified directional areas vary continuously, hence the visibility
bodies converge in homothetic distance. The colour separation margin
forces their selected ellipsoids to coincide eventually.
\end{proof}

\begin{lemma}[Bisecting a ball; CV, Lemma 11]
\label{lem:cv-isoperimetric}
Let a nonzero polynomial divide the unit ball into sign regions of
relative volumes $a,b$, where $a+b=1$. Then
\[
 \cH^{n-1}(Z_p\cap B)
 \geq\frac12\left(a^{(n-1)/n}+b^{(n-1)/n}-1\right)
                      \cH^{n-1}(S^{n-1}).
\]
In particular, a surface bisecting a ball has a dimensional positive
lower bound for its area in that ball.
\end{lemma}
\begin{proof}
Apply the isoperimetric inequality separately to the two sign regions.
Their boundaries are contained in the union of the outer sphere and
two copies of the separating surface. Subtract the sphere area and
divide by two. The non-strict inequality includes the cases $a=0$
and $b=0$; the printed strict inequality needs those degenerate cases
excluded.
\end{proof}

\section{The subunit critical exponent}

\begin{theorem}[The Wolff floor used in the sticky argument]
\label{thm:hairbrush}
\lean{Kakeya.Assouad.PureWZ2WolffVolumeFloor,Kakeya.Assouad.pure_wz2_direct_wolff_floor}
\leanok
\uses{def:all-scale-cwa,def:shading}
For every $\eps>0$ there are $\eta>0$ and $\delta_0>0$ such that
a nonempty family of $\delta$-tubes satisfying all-scale CWA with
error $\delta^{-\eta}$, with $\delta^\eta$-dense shading, obeys
\[
                         |E_Y|\geq\delta^{1/2+\eps}
                         \qquad(0<\delta\leq\delta_0).
\]
\end{theorem}
\begin{proof}
\leanok
This is the form of Wolff's $5/2$-dimensional hairbrush bound used
in WZ2 Section 6. The low-multiplicity case follows by integrating
the shading multiplicity. In the high-multiplicity case, select a
stem tube and the tubes meeting it. Angular decomposition separates
the resulting hairbrush into nearly planar pieces; planar overlap
bounds and the Wolff nonconcentration conditions give the square-root
gain. Density pruning and the all-scale cover absorb fixed and
logarithmic losses into $\eps$.
\end{proof}

\begin{remark}[Scope of the hairbrush statement]
\label{rem:hairbrush}
The displayed theorem records the closed consequence needed here,
not the full maximal-function theorem of Wolff~\cite{Wolff95}.
That article is not needed as an additional paid source for the
blueprint. If an admissible exponent were greater than $1/2$, the
floor with a smaller positive loss would contradict its defining
upper bound at sufficiently small scales. The proof uses a convenient
ceiling $3/4<1$ and the floor with loss $1/8$; sharpness of the
ceiling is unnecessary.
\end{remark}
